\section{Assumption grades}

\subsection{Syntax} 
 The extended syntax is given in \cref{fig:ext-syntax}. 
\begin{figure}
\begin{grammatica}
\produzione{\e}{\x\mid\FieldAccess{\x}{\f}\mid
\FieldUpdate{\x}{\f}{\y}\mid\Free{\x}\mid\MethCall{\x}{\m}{\xs}\mid\New{\C}{\xs}\mid\Let{\x}{\e}{\e'}}{expression}\\
\produzione{\tau}{\C}{non-graded type}\\
\produzione{\grdec}{\GrDec{\Rgr}{\rgr}{\varphi}}{grade decoration}\\
\produzione{\T}{\Graded{\tau}{\grdec}}{(graded) type}\\
\produzione{\varphi}{\FieldGrade{\f_1}{\Rgr_1}{\rgr_1}\ldots\FieldGrade{\f_n}{\Rgr_n}{\rgr_n}}{field grades}\
\end{grammatica}
\caption{Extended syntax}\label{fig:ext-syntax}
\end{figure}

In addition to the previous constructs, we include field update and deallocation.  

\subsection{Type system}

Information extracted from the class table and contexts are as before, apart that usage grades are replaced by pairs of assumption and usage grades. The sum of such pairs is defined as follows:

\begin{quote}
$\Pair{\Rgr_1}{\rgr_1}\rsum\Pair{\Rgr_2}{\rgr_2}=\Pair{(\Rgr_1 \rsub \rgr_2) \rmeet (\Rgr_2 \rsub \rgr_1)}{(\rgr_1\rsum\rgr_2)}$
\end{quote}

Such sum is a key technical feature of our approach. 
 Usually, in graded type systems, usage grades of occurrences of the same variable are summed. 

Here, assumption grades should be taken into account as well. Assume to have $\CtxElem{\x}{\Rgr_1}{\rgr_1}{\tau}$ and $\CtxElem{\x}{\Rgr_2}{\rgr_2}{\tau}$ in two contexts, respectively. In the first context, the assumption $\Rgr_1$ means that $\x$ can be used elsewhere with grade (at most) $\Rgr_1$.  When summing the contexts, $\x$ is used with grade $\rgr_2$, hence this amount should be subtracted from $\Rgr_1$. The same holds conversely. Altogether, $\x$ can still be used elsewhere with a grade which is the meet of the two residual assumptions. 
For a concrete example, consider the algebra of extended natural numbers. Assuming to have $\CtxElem{\x}{3}{2}{\tau}$ and $\CtxElem{\x}{2}{3}{\tau}$ in the two contexts, respectively, we get $3-3=0$ and $2-2=0$, hence we have $\CtxElem{\x}{0}{5}{}$ in the sum.  In \cref{ex:ex1} later we show an example of code where this situation arises. 

Note that, in this way, the sum is undefined when, for some variable, the assumption in one context is not enough to cover the usage in the other. This is an additional source of partiality, besides the case in which the types in the two contexts are different. 


The extended typing rules are given in \cref{fig:ext-typing-exp}. They use the following definitions:

\begin{figure}
\begin{small}
\begin{math}
\begin{array}{l}

%R';r'+S;s <= R;r

\NamedRule{t-var}{}
{ \TypingExp{\Gamma,\AUCtxElem{\x}{\Rgr}{\rgr}{\varphi}{\tau}}
{\x}
{\Graded{\tau}{\GrDec{\Sgr}{\sgr}{\varphi}}}
{\Gamma,\AUCtxElem{\x}{\Rgr'}{\rgr'}{\varphi}{\tau}}
}
{\Pair{\Rgr'}{\rgr'}\rsum\Pair{\Sgr}{\sgr}\rord\Pair{\Rgr}{\rgr}}
%\\[4ex]
%\NamedRule{t-field-access}{}
%{
%\TypingExp{\Gamma,\UCtxElem{\x}{\rgr}{\varphi,\FieldUGrade{\f_i}{\rgr_i}}{\C}}
%{\FieldAccess{\x}{\f_i}}
%{\Graded{\tau_i}{\UGrDec{\rgr\rmul\sgr}{\varphi_i}}}
%{\Gamma,\UCtxElem{\x}{\rgr}{\varphi,\FieldUGrade{\f_i}{\rgr'_i}}{\tau_i}}}
%{\fields{\C}=\UField{\tau_1}{\rgr_1}{\f_1} \ldots \UField{\tau_n}{\rgr_n}{\f_n}\\
%i\in 1..n\\
%\rgr'_i\rsum\sgr\rord\rgr_i\\
%\fgrades{\tau_i}=\varphi_i
%}
%\\[6ex]
%\mbox{\EZComm{regola del metodo solo abbozzata!}}
%\\
%\NamedRule{t-invk}
%{}{\TypingExp{\Gamma,\Gamma^\textit{pre}}{\MethCall{\x}{\m}{\x_1,\ldots,\x_n}}{\T}{\Gamma,\Gamma^\textit{post}}}{
%\Gamma^\textit{pre}=\CtxElem{\x}{\grdec}{\C},\CtxElem{\x_1}{\grdec_1}{\tau_1},\ldots,\CtxElem{\x_n}{\grdec_n}{\tau_n}\\
%\mtype{\C}{\m}=\mType{\Gamma_\m^\textit{pre}}{\T}{\Gamma_\m^\textit{post}}\\
%\Gamma_\m^\textit{pre}=\CtxElem{\this}{\grdec'}{\C},\CtxElem{\y_1}{\grdec'_1}{\tau_1},\ldots,\CtxElem{\y_n}{\grdec'_n}{\tau_n}\\
%\Gamma_\m^\textit{post}=\CtxElem{\this}{\grdec''}{\C},\CtxElem{\y_1}{\grdec''_1}{\tau_1},\ldots,\CtxElem{\y_n}{\grdec''_n}{\tau_n}\\
%\Gamma^\textit{post}=\ldots
%}
%\\[6ex]
%\NamedRule{t-new}{}
%{\TypingExp{\Gamma,\Gamma^\textit{pre}}
%{\New{\C}{\x_1,\ldots,\x_n}}
%{\UGraded{\C}{\rgr}{\varphi}}
%{\Gamma,\Gamma^\textit{post}
%}
%}
%{
%\fields{\C}=\UField{\tau_1}{\sgr_1}{\f_1} \ldots \UField{\tau_n}{\sgr_n}{\f_n}\\
%\Gamma^\textit{pre}=\UCtxElem{\x_1}{\rgr_1}{\varphi_1}{\tau_1},\ldots,\UCtxElem{\x_n}{\rgr_n}{\varphi_n}{\tau_n}\\
%\Gamma^\textit{post}=\UCtxElem{\x_1}{\rgr'_1}{\varphi_1}{\tau_1},\ldots,\UCtxElem{\x_n}{\rgr'_n}{\varphi_n}{\tau_n}\\
%}
%\\[5ex]
%\newcommand{\LocalGamma}{\hat{\Gamma}}
%\NamedRule{t-let}{
%\TypingExp{\Gamma}{\e}{\Graded{\tau}{\grdec}}{\Gamma'}\BigSpace\TypingExp{\Gamma',\LocalGamma}{\e'}{\T}{\Gamma''}
%}
%{\TypingExp{\Gamma}{\Let{\x}{\e}{\e'}}{\T}{\Gamma''}}
%{\LocalGamma\ctxord\CtxElem{\x}{\grdec}{\tau}
%%\\\T'\subt\Graded{\tau}{\Rgr}{\sgr}
%}
%
\end{array}
\end{math}
\end{small}
\caption{Extended typing rules for expressions}
\label{fig:ext-typing-exp} 
\end{figure}

\EZComm{CONTROLLARE SUBTYPING E WEAKENING}

In rule \refToRule{t-var}, besides consuming the usage grade as in the previous version, 

in addition to the behaviour on usage grades he availability of the resource denoted by $\x$ is $\rgr$ before evaluating the variable occurrence. 
Some amount $\sgr$ of the resource is used to obtain the result, letting a residual availability $\rgr'$. Of course the initial availability should be enough to allow the usage. 

\subsection{Examples} 

